\documentclass[palette=navy,titlestyle=rigorous,titlealign=center]{inkpaper}
% Other looks: titlestyle=open|panel, titlealign=left, palette=teal|slate

\usepackage{pgfplots}
\pgfplotsset{compat=1.18}

\title{Calibrated Critics: Learning When to Trust\\Rubric Rewards in Open-Ended RL}
\shorttitle{Calibrated Critics}
\papertag{Preprint}
\author{Alice Chen\textsuperscript{1,*}, Bo Li\textsuperscript{1,*},
  Carlos Diaz\textsuperscript{2}, Dana Park\textsuperscript{1,\dag}}
\affiliations{\textsuperscript{1}Affiliation One \quad
  \textsuperscript{2}Affiliation Two}
\authornotes{\textsuperscript{*}Equal contribution \quad
  \textsuperscript{\dag}Corresponding author}
\date{October 15, 2026}
\metadata{Version}{v1 (under review)}   % any custom label works
\correspondence{\href{mailto:author@example.com}{author@example.com}}
\code{https://github.com/your-org/calibrated-critics}
\model{https://huggingface.co/your-org}
\data{https://huggingface.co/datasets/your-org}
\project{https://your-org.github.io/calibrated-critics}

% Replace these text marks with institutional logo files if available:
% \leftlogo{\includegraphics[height=16pt]{logo/university.pdf}}; or use the
% built-in badge helper, \leftlogo{\placeholderlogo{U}{University Name}}.
\leftlogo{{\sffamily\bfseries\small\color{ipAccent}University Name}}
\rightlogo{\begin{tabular}{r}
  \sffamily\bfseries\scriptsize Research Lab\\[-1pt]
  \sffamily\scriptsize\color{ipGray}Institute Name
\end{tabular}}
\firstpagenote{Preprint.}

\abstract{Rubric-based reinforcement learning has become a standard way to
post-train language models on tasks without verifiable answers, yet the judge
that scores each rubric item is itself an imperfect proxy. We show that the
judge's confidence is informative: rubric items on which the judge is
uncertain account for most of the reward that a policy later exploits. We
introduce \emph{Calibrated Critics}, which reweights each criterion by an
estimate of the judge's reliability on that criterion and prompt. Across
long-form research and medical benchmarks, our method improves held-out gold
judge scores by up to 6.1 points while matching the baseline's in-domain
reward, and it substantially narrows the gap between training and gold
judges throughout training.}

\begin{document}
\maketitle

% Teaser: [h] keeps it under the title block; a [t] float would go above it.
\begin{figure}[h]
\centering
\begin{tikzpicture}
\begin{axis}[inkpaper bars,width=\linewidth,height=4.6cm,ymin=20,ymax=68,
  symbolic x coords={HealthBench-Hard,ResearchQA,ScholarQA,DeepResearch},
  xtick=data,ytick={20,40,60}]
\addplot coordinates
  {(HealthBench-Hard,31.2)(ResearchQA,48.9)(ScholarQA,52.4)(DeepResearch,38.0)};
\addplot coordinates
  {(HealthBench-Hard,36.8)(ResearchQA,53.1)(ScholarQA,57.0)(DeepResearch,41.5)};
\addplot coordinates
  {(HealthBench-Hard,38.0)(ResearchQA,55.7)(ScholarQA,57.9)(DeepResearch,43.2)};
\addplot+[nodes near coords] coordinates
  {(HealthBench-Hard,40.3)(ResearchQA,59.2)(ScholarQA,60.1)(DeepResearch,47.8)};
\legend{Base model,Uniform rubric,Rubric dropout,\textbf{Calibrated Critics}}
\end{axis}
\end{tikzpicture}
\caption{\captitle{Out-of-distribution results.} Gold judge scores on four
long-form benchmarks. Calibrated Critics improves over every baseline on all
four.}
\label{fig:teaser}
\end{figure}

\section{Introduction}

Reinforcement learning with verifiable rewards works well when a ground truth
exists~\citep{shao2024deepseekmath}. Many of the tasks we care most about,
such as explaining a research area or giving medical advice, are open-ended.
A common recipe is to write a rubric for each prompt and let an LLM judge
grade each criterion~\citep{gunjal2025rubrics}. The rubric, however, is a
proxy, and a policy optimized against it long enough learns to exploit the
difference~\citep{gao2023scaling}.

Our contributions are as follows:
\begin{itemize}
  \item We show that judge uncertainty predicts which rubric items will be
    exploited, before exploitation happens.
  \item We propose a per-criterion reliability weight that is cheap to
    estimate and drops into any GRPO-style objective.
  \item We evaluate on four long-form benchmarks and two model scales.
\end{itemize}

\begin{takeaway}[label=tk:main]
\textbf{Judge uncertainty identifies the rubric criteria most vulnerable to
reward hacking.} Most exploitation concentrates on a small set of criteria
where the judge was already uncertain before policy optimization.
\end{takeaway}

\section{Method}

\subsection{Reliability-weighted rubric reward}

Given a prompt $x$ with rubric $\{c_k\}_{k=1}^{K}$ and weights $w_k$, the
standard reward for a response $y$ is the weighted fraction of satisfied
criteria. We instead use
\begin{equation}
  r(x, y) \;=\; \frac{\sum_{k=1}^{K} w_k\,\rho_k(x)\,\mathbb{I}[\,y \models c_k\,]}
                    {\sum_{k=1}^{K} w_k\,\rho_k(x)},
\end{equation}
where $\rho_k(x) \in [0,1]$ is the estimated reliability of the judge on
criterion $c_k$.

\begin{observation}
\textbf{Judge disagreement is concentrated.} A small fraction of criteria
accounts for most of the judge's vote entropy, and these are predominantly
stylistic rather than factual criteria.
\end{observation}

\paragraph{Estimating reliability} We sample the judge $m$ times per
criterion and set $\rho_k$ to one minus the normalized entropy of its votes.

\begin{discussion}
\textbf{Relationship to rubric dropout.} Dropout makes a fixed rubric harder
to exploit by varying which criteria receive reward. Reliability weighting
addresses a complementary failure mode: it reduces the influence of criteria
whose judgments are unstable even before the policy begins to optimize
against them.
\end{discussion}

\begin{figure}[t]
\centering
\begin{tikzpicture}
\begin{axis}[width=0.78\linewidth,height=5.0cm,
  axis lines=left,xlabel={Training step},ylabel={Gold judge score},
  every axis label/.style={font=\sffamily\footnotesize,text=ipGray},
  tick label style={font=\sffamily\scriptsize,text=ipGray,
    /pgf/number format/assume math mode=true},
  axis line style={ipRule},major tick style={ipRule},
  ymajorgrids,grid style={ipRule!50},
  legend style={draw=none,font=\sffamily\scriptsize,fill=none,at={(0.02,0.98)},anchor=north west},
  legend cell align=left,xmin=0,xmax=500,ymin=40,ymax=64]
\addplot[ipAccent,line width=1.3pt,smooth] coordinates
  {(0,44)(100,51)(200,56)(300,59)(400,61)(500,62)};
\addlegendentry{Calibrated Critics}
\addplot[ipFaint,line width=1.1pt,smooth,dashed] coordinates
  {(0,44)(100,50.5)(200,54)(300,54.5)(400,53)(500,51.5)};
\addlegendentry{Uniform rubric}
\end{axis}
\end{tikzpicture}
\caption{\captitle{Training dynamics.} Gold judge score on held-out prompts
during training. The uniform rubric baseline peaks and then declines as the
policy exploits unreliable criteria; reliability weighting keeps improving.}
\label{fig:curve}
\end{figure}

\section{Experiments}

\paragraph{Setup} We post-train two policy scales with GRPO against three
baselines: the base model, the uniform rubric reward, and rubric dropout,
which masks a random subset of criteria at every step. We evaluate on four
held-out benchmarks:\footnote{All benchmarks are scored by a gold judge that
is never used during training; Appendix~\ref{app:judge} gives its
configuration.}
\begin{enumerate}
  \item \textbf{HealthBench-Hard}: open-ended clinical advice graded against
    physician-written rubrics.
  \item \textbf{ResearchQA} and \textbf{ScholarQA}: literature-synthesis
    questions with citation-level criteria.
  \item \textbf{DeepResearch}: multi-document report writing.
\end{enumerate}

\begin{table}[t]
\centering
\begin{threeparttable}
\caption{\captitle{Main results.} Gold judge scores on out-of-distribution
benchmarks. Best values are shown in bold.}
\label{tab:main}
\small
\begin{tabular}{lccccc}
\toprule
Method & HealthBench-Hard & ResearchQA & ScholarQA & DeepResearch & Avg. \\
\midrule
Base model          & 31.2 & 48.9 & 52.4 & 38.0 & 42.6 \\
Uniform rubric      & 36.8 & 53.1 & 57.0 & 41.5 & 47.1 \\
Rubric dropout      & 38.0 & 55.7 & 57.9 & 43.2 & 48.7 \\
\oursrow \textbf{Calibrated Critics} & \textbf{40.3} & \textbf{59.2} & \textbf{60.1} & \textbf{47.8} & \textbf{51.9} \\
\bottomrule
\end{tabular}
\end{threeparttable}
\end{table}

\subsection{Main results}
Table~\ref{tab:main} reports gold judge scores on all four benchmarks and
Figure~\ref{fig:curve} shows the training dynamics. The gain is largest on
DeepResearch, whose rubrics contain the most criteria the judge is uncertain
about, and smallest on ScholarQA, whose criteria are mostly citation checks
that the judge grades consistently.

\begin{finding}
Reliability weighting improves gold judge scores on every benchmark while
leaving in-domain rubric reward unchanged, consistent with
Takeaway~\ref{tk:main}.
\end{finding}

\subsection{How many judge samples are needed?}
Figure~\ref{fig:ablation} varies the number of judge samples $m$ used to
estimate $\rho_k$. Most of the benefit arrives by $m=3$; beyond $m=5$ the
estimate is stable and the extra judge calls are wasted.

% [htb]: in place after the paragraph that cites it when it fits, else the
% next page top. The inkpaper preset supplies fonts, axis colours and grid;
% per-series keys (marks, dashes) and the axis limits stay explicit.
\begin{figure}[htb]
\centering
\begin{tikzpicture}
\begin{axis}[inkpaper,width=0.6\linewidth,height=3.6cm,
  xlabel={Judge samples $m$},ylabel={Avg.\ gold score},
  xtick={1,3,5,9},xmin=1,xmax=9,ymin=46,ymax=53]
\addplot[ipAccent,line width=1.3pt,mark=*,mark size=1.8pt] coordinates
  {(1,47.4)(3,51.0)(5,51.9)(9,52.0)};
\addplot[ipFaint,dashed,line width=1pt,domain=1:9] {47.1};
\end{axis}
\end{tikzpicture}
\caption{\captitle{Judge-sample ablation.} Average gold score as a function of
$m$; the dashed line is the uniform rubric baseline.}
\label{fig:ablation}
\end{figure}

\begin{practicalguidance}
\textbf{Calibrate the judge before starting reinforcement learning.} Sample
each rubric criterion several times on representative prompts, estimate its
agreement, and downweight unstable criteria. Recompute these estimates
periodically as the policy distribution shifts.
\end{practicalguidance}

\section{Limitations}
Reliability weighting assumes that judge disagreement signals a weak criterion
rather than a hard response; when the two coincide, it down-weights exactly
the criteria that matter most.

\bibliographystyle{plainnat}
\bibliography{references}

\beginappendix

\section{Implementation details}

\subsection{Training hyperparameters}
We train with GRPO using 8 rollouts per prompt, a learning rate of
$10^{-6}$, and a KL coefficient of 0.001 for 500 steps.

\subsection{Judge configuration}
\label{app:judge}
The training judge is sampled $m=5$ times per criterion at temperature 0.7.

\section{Prompts}

\begin{promptbox}[Rubric judge prompt]
You are grading a response against a single rubric criterion.\\[0.4em]
Question: \{question\}\\
Response: \{response\}\\
Criterion: \{criterion\}\\[0.4em]
Answer YES if the response satisfies the criterion, otherwise NO.
\end{promptbox}

The judge's output format is fixed with a second, machine-readable prompt,
pasted verbatim:
\begin{promptbox*}[Judge output format]
Question: {question}
Response: {response}
Criterion: {criterion}

Return JSON: {"verdict": "YES" | "NO", "confidence": 0-100}
# confidence is a percentage; use max_tokens=64.
\end{promptbox*}

\section{Additional results}
Per-benchmark training curves are shown in the supplementary material.

\end{document}
